%%%%

%%%   Самарский государственный технический университет

%%%   Кафедра прикладной математики и информатики

%%%

%%%   Преамбула для статьи в журнал "Вестник Самарского государственного

%%%   технического университета. Серия: «Физико-математические науки»"

%%%   Ориентирована под MikTEx 2.4 for Win32

%%%

%%%   Хотя сам журнал собирается под GNU/Linux на дистрибутиве TeXLive



\documentclass[11pt,a4paper,twoside]{article}



\usepackage[cp1251]{inputenc}

\usepackage[T2A]{fontenc}

\usepackage[english,russian]{babel}

\usepackage{samgtubib}

\usepackage[dvips]{graphicx}

\usepackage[multiple]{footmisc}



\usepackage{samgtu}

\renewcommand{\VestnikYear}{2009} % Год издания Вестника

\renewcommand{\VestnikNumber}{2\,(19)} % Номер Вестника

\renewcommand{\VestnikMonth}{Ноябрь} % Месяц выхода Вестника

\renewcommand{\VestnikData}{01.11.09} % Дата подписания Вестника в печать

\renewcommand{\VestnikUPL}{26,64} % Усл. печ. л.

\renewcommand{\VestnikUIL}{25,73} % Уч.-изд. л.

\renewcommand{\VestnikRegNumber}{479/09} % Регистрационный номер

\renewcommand{\VestnikOrder}{994} % Номер заказа






%
%\begin{document}


\renewcommand{\firstpage}{92}
\renewcommand{\lastpage}{109}



\udk{539.376}
\msc{74K20, 74R20}

\title{Моделирование установившейся ползучести пространст\-венно"=армированных металлокомпозитов с~уч\"eтом анизотропии фазовых материалов}{Моделирование установившейся ползучести  \dots}

\author{А.\,П.~Янковский}{Я\,н\,к\,о\,в\,с\,к\,и\,й~А.\,П.}

\institute{Институт теоретической и прикладной механики им. С.\,А.~Христиановича СО РАН,\\
630090, Россия, Новосибирск, ул.~Институтская, 4/1.}


\email{E-mail}{lab4nemir@rambler.ru}

\otherinf{\fullauthor{Андрей Петрович Янковский}\regalia{д.ф.-м.н., проф.} \position{ведущий научный сотрудник}
\dept{лаб. физики быстропротекающих процессов}}

\keywords{металлокомпозиты, армирование пространственное, установившаяся ползучесть, анизотропия, структурная теория}

\workabstract{Предложена итерационная модель, описывающая
механическое поведение металлокомпозитных сред, перекрёстно
армированных в пространстве, материалы компонентов которых
работают в условиях установившейся анизотропной ползучести.
Проведён сравнительный анализ расчётов, выполненных по разным
структурным моделям механического поведения перекрестно
армированных в плоскости и в пространстве металлокомпозитов в
условиях установившейся ползучести.}

\engtitle{Modeling of steady creep of 3D reinforced metal-composits with anisotropy of phase materials}

\engauthor{A.\,P.~Yankovskii}{Y\,a\,n\,k\,o\,v\,s\,k\,i\,i~A.\,P.}


\enginstitute{Khristianovich Institute of theoretical and applied
mechanics \\
the Siberian Branch of the Russian Academy of Science,\\
 4/1, Institutskaya st., Novosibirsk,  630090, Russia.}


\otherenginf{\fullauthor{Andrey P. Yankovskii} \regalia{Dr. Sci. (Phys. \& Math.)}
\position{Leading Research Scientist} \dept{Lab. of Fast Processes Physic}}

\engabstract{The iteration model describing mechanical behavior of cross-reinforced in the space metal-composites with steady-state anisotropic creep is presented. The comparative analysis of the calculations on different structural models of mechanical behavior of cross-reinforced in plane and in space metal-composites with steady-state creep is carried out.
}

\engkeywords{metal-composite, 3D reinforcement, steady creep, anisotropy, structural theory}

\date{25/X/2011}
\revisiondate{27/II/2012}

\maketitle


\Section[n]{Введение} В последнее время особое внимание к себе
привлекают композиционные материалы с пространственным
расположением арматуры [1,~2~и~др.]. 
Эффективность пространственного расположения арматуры определяется не только
возможностью ликвидировать такой недостаток слоистых композитов
как опасность расслоения вследствие слабого сопротивления
межслойному сдвигу и поперечному отрыву, но и возможностью
локализовать в пределах нескольких пространственных ячеек
распространение трещин. Этим резко повышается несущая способность
материала в толстостенных конструкциях, особенно в зонах
приложения локализованных нагрузок и концентраторов напряжений при
интенсивных термосиловых воздействиях, характерных для современных
технических устройств \cite{yank:bib1}. При наличии таких
воздействий конструкцию целесообразно изготавливать из
металлокомпозитов, которые находят все более широкое применение на
практике.

Известно, что при длительном термосиловом нагружении
металлоконструкции большую часть времени эксплуатации работают в
условиях установившейся ползучести \cite{yank:bib3}, поэтому
актуальной является проблема математического моделирования
процессов ползучести в армированных металлокомпозитах. На
сегодняшний день эта проблема находится, по сути, в зачаточном
состоянии. Так, в работе \cite{yank:bib4} рассмотрена
установившаяся ползучесть металлокомпозитов слоистой структуры; в
\cite{yank:bib5} предложена модель ползучести однонаправленно
армированного композита, но такие виды армирования редко
встречаются на практике (как правило, в стержневых элементах
конструкций); в \cite{yank:bib6} разработана модель механического
поведения в условиях установившейся ползучести перекрестно
армированного в плоскости металлокомпозитного слоя с одномерным
напряжённым состоянием в армирующих элементах, в рамках которой не
учитываются напряжения в поперечном направлении в арматуре (как
нормальные, так и касательные). Последнее обстоятельство как бы
моделирует неидеальное сцепление арматуры со связующим, что
нередко наблюдается на практике при перекрестном армировании.
Современные же технологии позволяют создавать композиционные
материалы с хорошей адгезией \cite{yank:bib5}, близкой к
идеальному механическому сцеплению.

Все указанные модели [4--6] базируются на предположении об
изотропии материалов компонентов композиции. Однако известно, что
при испытаниях на ползучесть анизотропия проявляется больше, чем
при любом другом виде механических испытаний \cite{yank:bib3}.
Анизотропия характеристик ползучести может быть не только связана 
с исходным состоянием материала, но и порождена кратковременными
пластическими деформациями, возникающими при нагружении и
предшествующими ползучести, а также может возникнуть на первой
стадии ползучести, сопутствуя упрочнению.

В связи с указанными обстоятельствами настоящее исследование посвящено
разработке модели механического поведения пространственно"=армированных
металлокомпозитов, работающих в условиях установившейся ползучести всех
фазовых материалов, которые предполагаются анизотропными и в которых
учитываются все компоненты напряжённого состояния.

\smallskip

\Section{Итерационная модель пространственно"=армированного
металлокомпозита, работающего в условиях установившейся
анизотропной ползучести} Так как наличие арматуры с существенно
разными механическими характеристиками значительно расширяет
диапазон свойств композиционных материалов с пространственной
схемой армирования \cite{yank:bib1}, в глобальной декартовой
системе координат ($x_1$, $x_2$, $x_3 $) рассмотрим гибридный
металлокомпозит, армированный в произвольных направлениях $K$
семействами прямолинейных волокон (проволок) с относительным
объёмным содержанием "--- плотностями (интенсивностями)
армирования $\omega _k $, $1\le k\le K$.
Относительное объёмное содержание связующего материала обозначим
через $\omega _0$, тогда имеет место условие нормировки
\begin{equation}
\label{yank:eq1}
\omega _0 +\sum\limits_k \omega_k
=1,\qquad \omega _0>0,\,\omega _k \ge 0,\; 1\le k\le K.
\end{equation}
Здесь и далее суммирование осуществляется по указанному индексу от 1 до $K$,
если не проставлены пределы.

Кроме условия нормировки \eqref{yank:eq1} должны выполняться и физические условия
взаимного непроникновения материалов различных компонентов композиции. Эти
условия накладывают определенные ограничения на предельно допустимые
значения суммарных плотностей армирования (на значения суммы в (\ref{yank:eq1})) при
плотной упаковке армирующих элементов. Так, в \cite{yank:bib1} приведены указанные
предельные значения для некоторых структур пространственного армирования
композитной среды, которые меньше единицы. Далее в настоящем исследовании
предполагается, что эти ограничения на значения суммарных плотностей
армирования выполняются. При построении модели механического поведения
рассматриваемого металлокомпозита знание конкретных чисел $\omega _k$
необязательно, важным является выполнение условия нормировки~(\ref{yank:eq1}).

\begin{wrapfigure}{O}{.45\textwidth}
\includegraphics[width=.4\textwidth]{./00PIC/yank1.eps}
\caption{Локальная система координат, связанная с арматурой $k$-того семейства}
\end{wrapfigure}

С каждым $k$-тым семейством арматуры свяжем свою локальную ортогональную систему
координат ($x_1^{(k)}$, $x_2^{(k)}$, $x_3^{(k)}$, $1\le k \le K$) так, чтобы ось $x_1^{(k)} $ совпадала с~направлением траекторий армирования этого семейства, а оси $x_2^{(k)}$, $x_3^{(k)}$ были перпендикулярны этим траекториям (рис.~1).
Направляющие косинусы (см.~ табл. 11.1 в \cite{yank:bib7}) между осями $x_i^{(k)}$ и $x_j$ обозначим как $l_{ij}^{(k)}$ ($i,\,j=1,\,2,\,3)$, для которых
выполняются равенства
\[
\sum\limits_{j=1}^3 l_{ij}^{(k)} l_{ij}^{(k)} = \sum\limits_{j=1}^3
l_{ji}^{(k)} l_{ji}^{(k)}=1,\] \[ i=1,\,2,\,3,\; 1\le k \le K.
\]




Все компоненты композиции предполагаются обладающими анизотропией
ползучести, причем каждый $k$-тый фазовый материал в условиях ползучести
характеризуется потенциалом ползучести $\Phi _k (S_k,\,\Theta _k)$, зависящим только от температуры $\Theta _k$ и от одного
квадратичного инварианта $S_k$, составленного из компонентов тензора
напряжений $\sigma _{ij}^{(k)} $ и компонентов некоторого тензора
анизотропии четвертого ранга $A_{ijlm}^{(k)} (\Theta _k)$,
также зависящего от температуры. Запишем этот инвариант следующим образом
(см.~стр.~283 в \cite{yank:bib3}):
\begin{equation}
\label{yank:eq2}
S_k =\sum\limits_{i=1}^3
\sum\limits_{j=1}^3
\sum\limits_{l=1}^3
\sum\limits_{m=1}^3 A_{ijlm}^{(k)} (\Theta _k)\sigma _{ij}^{(k)} \sigma _{lm}^{(k)},
\quad 0\le k\le K,
\end{equation}
где
\begin{equation}
\label{yank:eq3}
A_{ijlm}^{(k)} =A_{jilm}^{(k)} =A_{ijml}^{(k)}
=A_{lmij}^{(k)}.
\end{equation}

Скорости деформаций установившейся ползучести $\xi _{ij}^{(k)} $ в $k$-том
фазовом материале с учётом \eqref{yank:eq2}, \eqref{yank:eq3} определяются так:
\begin{equation}
\label{yank:eq4}
\xi _{ij}^{(k)} {=}\frac{\partial \Phi _k (S_k,\,\Theta _k)}{\partial S_k }\frac{\partial S_k
}{\partial \sigma _{ij}^{(k)}}=2\Phi'_k (S_k,\,\Theta _k)\sum\limits_{l=1}^3 \sum\limits_{m=1}^3
A_{ijlm}^{(k)} (\Theta _k)\sigma _{lm}^{(k)},\, i,\,j{=}1,\,2,\,3,
\end{equation}
где
\[
\Phi'_k (S_k,\,\Theta _k)\equiv \frac{\partial \Phi _k (S_k,\,\Theta _k)}{\partial S_k },\quad 0\le k\le
K.
\]

Если дополнительно предполагается, что деформации ползучести удовлетворяют
условию несжимаемости, то получим \cite{yank:bib3}, что
\begin{equation}
\label{yank:eq5}
\sum\limits_{i=1}^3 A_{iilm}^{(k)}
=\sum\limits_{i=1}^3 A_{iiml}^{(k)} = \sum\limits_{i=1}^3
A_{lmii}^{(k)} = \sum\limits_{i=1}^3 A_{mlii}^{(k)}
=0,\quad 0\le k\le K.
\end{equation}

Для удобства изложения (хотя это и не принципиально) характеристики
анизотропии связующего ($k=0$) считаются заданными в глобальной системе
координат ($x_1$, $x_2$, $x_3$), а арматуры $k$-того семейства "--- в локальной
системе ($x_1^{(k)}$, $x_2^{(k)}$, $x_3^{(k)}$, $1\le k\le K$).
Определяющие соотношения установившейся ползучести для металлокомпозиции
будем строить также в глобальной системе координат. В связи с этим далее
чертой сверху будем обозначать величины, определенные в локальной системе
координат ($x_1^{(k)}$, $x_2^{(k)}$, $x_3^{(k)}$), связанной с $k$-тым
семейством волокон, а без черты "--- определенные в глобальной системе ($x_1$, $x_2$, $x_3$) (т.\,е. соотношения
\eqref{yank:eq2}--\eqref{yank:eq5} записаны в глобальной
системе; если же компоненты тензоров $A_{ijlm}^{(k)}$, $\sigma _{lm}^{(k)}$, $\xi _{ij}^{(k)} $ снабдить сверху чертой, то получим запись соотношений
(\ref{yank:eq2})--(\ref{yank:eq5}) в $k$-той локальной системе, $1\le k\le K$).

Определяющие соотношения \eqref{yank:eq4} в указанных системах отсчёта можно записать в
матричной форме
\begin{equation}
\label{yank:eq6}
\xi _0 =2 \Phi'_0 (S_0,\,\Theta _0)A_0 \sigma _0,\quad \bar {\xi}_k =2 \Phi'_k (S_k,\,\Theta _k) \bar {A}_k  \bar{\sigma}_k,\quad 1\le k\le K,
\end{equation}
где
\begin{equation}
\label{yank:eq7}
\begin{array}{l}
\sigma_0^\ast =\left\{ \sigma _1^{(0)},\sigma _2^{(0)},\sigma
_3^{(0)},\sigma_4^{(0)},\sigma _5^{(0)},\sigma _6^{(0)}\right\}\equiv \left\{ \sigma _{11}^{(0)},\sigma _{22}^{(0)},\sigma _{33}^{(0)},\sigma _{23}^{(0)},\sigma _{31}^{(0)},
\sigma _{12}^{(0)}\right\}, \\
\xi_0^\ast =\left\{ \xi_1^{(0)},\xi _2^{(0)},\xi_3^{(0)},\xi _4^{(0)},\xi _5^{(0)},\xi _6^{(0)}\right\}\equiv
\left\{ \xi _{11}^{(0)},\xi _{22}^{(0)},\xi _{33}^{(0)},\xi_{23}^{(0)},\xi _{31}^{(0)},\xi _{12}^{(0)} \right\}, \\
\bar {\sigma }_k^\ast =\left\{\bar {\sigma }_1^{(k)},\bar {\sigma}_2^{(k)},\bar {\sigma }_3^{(k)},\bar {\sigma}_4^{(k)},\bar
{\sigma }_5^{(k)},\bar {\sigma }_6^{(k)} \right\}\equiv \left\{ \bar {\sigma }_{11}^{(k)},\bar {\sigma }_{22}^{(k)},\bar {\sigma
}_{33}^{(k)},\bar {\sigma }_{23}^{(k)},\bar {\sigma }_{31}^{(k)},\bar {\sigma }_{12}^{(k)} \right\}, \\
\bar {\xi }_k^\ast =\left\{\bar {\xi }_1^{(k)},\bar {\xi }_2^{(k)},\bar {\xi}_3^{(k)},\bar {\xi }_4^{(k)},\bar {\xi }_5^{(k)},
\bar {\xi }_6^{(k)} \right\}\equiv \left\{ \bar {\xi }_{11}^{(k)},\bar {\xi }_{22}^{(k)},\bar {\xi }_{33}^{(k)},\bar {\xi}_{23}^{(k)},\bar {\xi}_{31}^{(k)},
\bar {\xi }_{12}^{(k)}\right\};
\end{array}
\end{equation}
$A_0 =\bigl( a_{ij}^{(0) } \bigr)$, $\bar {A}_k =\bigl( \bar {a}_{ij}^{(k)} \bigr)$ "---
($6{\times}6$)"=матрицы (в общем случае несимметричные), элементы которых зависят
от температуры и определяются компонентами тензоров анизотропии ползучести
$A_{ijlm}^{(0)} (\Theta _0)$, $\bar {A}_{ijlm}^{(k)} (\Theta _k)$ (см.~(\ref{yank:eq4}), \eqref{yank:eq6}, \eqref{yank:eq7}):
\begin{equation}
\label{yank:eq8}
\begin{array}{llll}
a_{ij}^{(0)} =A_{iijj}^{(0)},   & 
a_{i4}^{(0)} = 2A_{ii23}^{(0)}, & 
a_{i5}^{(0)} = 2A_{ii31}^{(0)}, & 
a_{i6}^{(0)} = 2A_{ii12}^{(0)}, 
\\
a_{4i}^{(0)} = A_{23ii}^{(0)},  &
a_{5i}^{(0)} {=}A_{31ii}^{(0)}, &
a_{6i}^{(0)} =A_{12ii}^{(0)},   &
a_{44}^{(0)} =2A_{2323}^{(0)},  
\\
a_{45}^{(0)} =2A_{2331}^{(0)},  & 
a_{46}^{(0)}=2A_{2312}^{(0)},   &
a_{54}^{(0)} =2A_{3123}^{(0)}, &
a_{55}^{(0)} =2A_{3131}^{(0)}, 
\\
a_{56}^{(0)} =2A_{3112}^{(0)}, &
a_{64}^{(0)} =2A_{1223}^{(0)}, &
a_{65}^{(0)}=2A_{1231}^{(0)}, &
a_{66}^{(0)} =2A_{1212}^{(0)}, 
\\
\bar {a}_{ij}^{(k)} =\bar {A}_{iijj}^{(k)},  &
\bar {a}_{i4}^{(k)} =2\bar {A}_{ii23}^{(k)}, &
\bar {a}_{i5}^{(k)} =2\bar {A}_{ii31}^{(k)}, &
\bar {a}_{i6}^{(k)} =2\bar {A}_{ii12}^{(k)},
\\
\bar {a}_{4i}^{(k)} =\bar {A}_{23ii}^{(k)}, &
\bar {a}_{5i}^{(k)} =\bar {A}_{31ii}^{(k)}, &
\bar {a}_{6i}^{(k)} =\bar {A}_{12ii}^{(k)}, &
\bar {a}_{44}^{(k)} =2\bar {A}_{2323}^{(k)},
\\
\bar {a}_{45}^{(k)} =2\bar {A}_{2331}^{(k)}, &
\bar {a}_{46}^{(k)} =2\bar {A}_{2312}^{(k)}, &
\bar {a}_{54}^{(k)} =2\bar {A}_{3123}^{(k)}, &
\bar {a}_{55}^{(k)} =2\bar {A}_{3131}^{(k)}, 
\\
\bar {a}_{56}^{(k)} =2\bar {A}_{3112}^{(k)}, &
\bar {a}_{64}^{(k)} =2\bar {A}_{1223}^{(k)}, &
\bar {a}_{65}^{(k)} =2\bar {A}_{1231}^{(k)}, &
\bar {a}_{66}^{(k)} =2\bar {A}_{1212}^{(k)}, 
\\
 & i,\,j=1,\,2,\,3, & 1\le k\le K;
 \end{array}
\end{equation}
$\sigma _0$, $\bar {\sigma }_k $ "--- шестикомпонентные вектор"=столбцы напряжений в~связующем и арматуре $k$-того семейства; $\xi _0$, $\bar {\xi }_k $ "---
шестикомпонентные вектор"=столбцы скоростей деформаций установившейся
ползучести в тех же материалах соответственно (здесь в отличие от \cite{yank:bib7} для
удобства изложения под компонентами $\xi _4^{(0)}$, $\xi _5^{(0)}$, $\xi_6^{(0)}$ и $\bar {\xi }_4^{(k)}$, $\bar {\xi }_5^{(k)}$, $\bar {\xi
}_6^{(k)}$ векторов $\xi _0$, $\bar {\xi }_k $ понимаются не полные
сдвиговые скорости деформаций ползучести, а их половины, равные
соответствующим компонентам тензоров скоростей деформаций $\xi _{23}^{(0)}$, $\xi _{31}^{(0)}$, $\xi _{12}^{(0)}$ и $\bar {\xi }_{23}^{(k)}$,
$\bar {\xi }_{31}^{(k)}$, $\bar {\xi }_{12}^{(k)}$, что согласно (\ref{yank:eq4}),
(\ref{yank:eq6}) и приводит к несимметричности матриц $A_0$, $\bar {A}_k $ в случае
общей анизотропии, см.~(\ref{yank:eq8})); звездочка означает операцию транспонирования.

Равенства (\ref{yank:eq7}) задают соответствия между шестью компонентами $f_i^{(0)}$, $\bar {f}_i^{(k)}$ $(i=1, 2, \dots, 6)$
некоторых векторов ${\rm{\bf f}}_0$, ${\rm {\bf \bar {f}}}_k$ и компонентами соответствующих
симметричных тензоров второго ранга $f_{ij}^{(0)}$, $\bar {f}_{ij}^{(k)}$; $i,\,j=1,\,2,\,3$, $1\le k\le K$.

В реальных материалах наличие текстуры, созданной при обработке,
предполагает тот или иной вид симметрии свойств ползучести. Наиболее высокая
степень симметрии наблюдается в образцах из пруткового материала или
цилиндрических образцах, упрочненных путем растяжения вдоль оси \cite{yank:bib3}. Это
обстоятельство относится и к армирующим металлическим проволокам. Учитывая,
что ось волокна $k$-того семейства совпадает с направлением $x_1^{(k)}$,
ненулевые компоненты симметричных в этом случае матриц $\bar {A}_k$ при
выполнении условий несжимаемости (\ref{yank:eq5}) имеют выражения (см.~(69.4) в \cite{yank:bib3}):
\begin{equation}
\label{yank:eq9}
\begin{array}{ll}
\bar {a}_{11}^{(k)} = {2}/{3}, & 
\bar {a}_{1i}^{(k)} =\bar {a}_{i1}^{(k)} =-{1}/{3}, 
\\
\bar {a}_{ii}^{(k)} = ({ 2-\alpha _k})/{3}, & \bar {a}_{23}^{(k)} =\bar{a}_{32}^{(k)} = ({\alpha _k -1})/{3}, 
\\
\bar {a}_{44}^{(k)} = ({3/2-\alpha _k})/{3}, &
\bar {a}_{55}^{(k)} =\bar {a}_{66}^{(k)}=({3/2-\beta _k })/{3}, \\
\hfill i=2,\,3,  & 1\le k\le K,
\end{array}
\end{equation}
где $\alpha _k =\alpha _k (\Theta _k)$, $\beta _k =\beta _k (\Theta _k)$ "--- величины, характеризующие анизотропию
ползучести монотропных (трансверсально"=изотропных) волокон. 
Случай изотропных волокон получается из \eqref{yank:eq9}, если положить $\alpha _k =\beta _k
=0$. Требование положительной определенности квадратичных форм $S_k$ (см.~(\ref{yank:eq2})) с учётом \eqref{yank:eq8}, \eqref{yank:eq9} приводит к следующим ограничениям
\cite{yank:bib3}: $\alpha_k<3/2$, $\beta _k <3/2$, $1\le k\le K$.

В случае монотропных волокон, когда анизотропия ползучести в них определяется соотношениями (\ref{yank:eq9}), направляющие косинусы $l_{ij}^{(k)}$ можно однозначно определить с помощью двух углов сферической системы координат (см.~рис.~1): полярного расстояния $\theta _k $ и долготы $\varphi _k$. 
При этом ось $x_2^{(k)}$ удобно получить поворотом оси $x_2$ на угол $\varphi _k$ вокруг оси $x_3$ (именно этот случай изображен на рис.~1), а направление оси $x_3^{(k)}$ определяется векторным произведением ортов, задающих направления $x_1^{(k)}$, $x_2^{(k)}$. Направляющие косинусы $l_{ij}^{(k)}$ при таком задании локальной системы координат вычисляются по
формулам
\begin{equation}
\label{yank:eq10}
\begin{array}{ll}
l_{11}^{(k)} =\sin \theta _k \cos \varphi _k , &
l_{12}^{(k)} =\sin \theta _k \sin \varphi _k ,\\
l_{13}^{(k)} =\cos \theta _k , & 
l_{21}^{(k)} =-\sin \varphi _k ,\\
l_{22}^{(k)} =\cos \varphi _k, &
l_{23}^{(k)} =0,\\
l_{31}^{(k)} =-\cos \theta _k \cos \varphi _k, &
l_{32}^{(k)} =-\cos \theta _k \sin \varphi _k ,\\
l_{33}^{(k)} =\sin \theta _k ,  & 1\le k\le K.
\end{array}
\end{equation}

Соотношения \eqref{yank:eq10} могут быть использованы и в случае общей анизотропии
материалов арматуры, но тогда все компоненты тензора анизотропии ползучести
$A_{ijlm}^{(k)}$ волокон $k$-того семейства обязательно должны быть заданы именно в этой системе координат.

Так как установить фактическое распределение напряжений и скоростей
деформаций ползучести в композитной среде, где основной материал имеет
многочисленные более жесткие включения, весьма затруднительно, то при
построении практически пригодных определяющих уравнений установившейся
ползучести рассматриваемого металлокомпозита необходимо сделать некоторые
допущения в виде исходных предпосылок.

\begin{itemize}


\item[1.] Армированный материал представляет собой сплошное
макроскопически квазиоднородное анизотропное тело. (При достаточно
густом равномерном насыщении связующего арматурными стержнями или
волокнами это предположение вполне допустимо. К этому выводу
приходят все исследователи, изучающие механические свойства
дисперсно"=армированных сред \cite{yank:bib7}.)

\item[2.] Между связующим и арматурой существует полное сцепление "---
идеальный термомеханический контакт.

\item[3.] В пределах представительного элемента, выделенного из композита
на миниуровне, скорости деформаций ползучести, напряжения и
температуры во всех компонентах и в композиции кусочно"=постоянны.
Эффектами высших порядков, связанными с изменением полей скоростей
деформаций, напряжений и температур на микроуровне в малых
окрестностях границ контакта связующего и волокон, пренебрегаем.

\item[4.] Усредненные поля скоростей деформаций ползучести, напряжений и
температуры в композиции определяются по правилу простой смеси
"--- пропорционально объёмному содержанию $\omega _k $ каждого
составляющего.

\item[5.] Все фазовые материалы однородны, а установившаяся анизотропная
ползучесть в них описывается квазилинейными уравнениями
(\ref{yank:eq6}) с учётом (\ref{yank:eq7})--(\ref{yank:eq9}),
(\ref{yank:eq2}), (\ref{yank:eq3}).

\item[6.] Соотношения (\ref{yank:eq6}) с учётом
(\ref{yank:eq7})--(\ref{yank:eq9}), (\ref{yank:eq2}),
(\ref{yank:eq3}) удовлетворяют достаточным условиям сходимости
метода последовательных приближений, аналогичного методу секущего
модуля \cite{yank:bib8}.

\item[7.] К рассматриваемому моменту времени деформации ползучести
получили настолько значительное развитие, что по сравнению с ними
можно пренебречь мгновенными упругими и пластическими деформациями
\cite{yank:bib3,yank:bib8}.

\end{itemize}

Так как при решении пространственных задач установившейся
ползучести целесообразно использовать итерационный процесс,
позволяющий линеаризовать определяющие соотношения
(\ref{yank:eq6}) (или, что то же самое, (\ref{yank:eq4})) и
являющийся, по сути, обобщением метода секущего модуля, широко
используемого в случае изотропной ползучести \cite{yank:bib8}, 
в силу допущения 6 все дальнейшие рассуждения относятся к
реализации этого метода. В случае изотропных фазовых материалов
достаточные условия сходимости такого итерационного процесса
известны (см. стр. 199 в \cite{yank:bib9}); в случае же
анизотропной ползучести эти условия сходимости автору неизвестны,
но на практике проверку сходимости такого метода можно
контролировать в процессе решения конкретной задачи об
установившейся ползучести пространственно"=армированного
металлокомпозита.

Согласно допущению 6 будем предполагать, что на некоторой $n$-ной
итерации приближения для напряжений $\mathop {\sigma _{ij}^{(0)}
}\limits^n $, $\mathop {\bar {\sigma }_{ij}^{(k)} }\limits^n $ и
температуры $\Theta _0 $, $\Theta _k $ во всех компонентах
композиции уже известны, т.\,е. в силу соотношений типа
(\ref{yank:eq2}) известны приближения квадратичных форм $\mathop
{S_0 }\limits^n $, $\mathop {S_k }\limits^n $, $1\le k\le K$, а
значит, следующие ($n+1)$-е приближения неизвестных векторных
функций $\mathop {\xi _0 }\limits^{n+1} $, $\mathop {\bar {\xi }_k
}\limits^{n+1} $ и $\mathop {\sigma _0 }\limits^{n+1} $, $\mathop
{\bar {\sigma }_k }\limits^{n+1} $ удовлетворяют линейным
уравнениям (см.~(\ref{yank:eq6})):
\begin{equation}
\label{yank:eq11} \mathop {\xi _0 }\limits^{n+1} =2 \Phi'_0 \Bigl(
\mathop {S_0}\limits^n,\,\Theta _0 \Bigr)A_0 \mathop {\sigma _0
}\limits^{n+1},\,\mathop {\bar {\xi }_k }\limits^{n+1} =2 \Phi'_k
\Bigl( \mathop {S_k }\limits^n ,\,\Theta _k \Bigr) \bar {A}_k
\mathop {\bar {\sigma }_k }\limits^{n+1},\quad 1\le k\le K.
\end{equation}

В силу соотношений \eqref{yank:eq11} и допущений 1, 6 определяющие
уравнения установившейся ползучести для рассматриваемого
металлокомпозита на ($n+1)$-й итерации также можно записать в
матричной форме:
\begin{equation}
\label{yank:eq12}
\mathop \xi \limits^{n+1} =\mathop A\limits^n
\mathop \sigma \limits^{n+1} ,
\end{equation}
где $\sigma $, $\xi $ "--- шестикомпонентные вектор"=столбцы
усредненных напряжений и скоростей деформаций установившейся
ползучести в композиции, аналогичные $\sigma _0 $, $\xi _0 $ в
(\ref{yank:eq7}), где нужно отбросить индекс <<0>>; $\mathop
A\limits^n =\bigl( \mathop {a_{ij}}\limits^n \bigr)$ "--- $6
\times 6$-матрица, компоненты которой подлежат определению на
основе известных компонентов матриц $2 \Phi'_0 \Bigl(\mathop {S_0 }\limits^n ,\,\Theta _0 \Bigr)A_0 $, 
$2 \Phi'_k \Bigl(\mathop{S_k }\limits^n,\,\Theta _k \Bigr)\bar {A}_k,$ $1\le k\le K$.

Далее в промежуточных формулах для сокращения записи будем
опускать верхние индексы $n$ и $n+1$, обозначающие номер итерации,
памятуя о том, что в (\ref{yank:eq11}) значения квадратичных форм
$S_0 $, $S_k $, а значит и функций $\Phi'_0$,~$\Phi'_k $, $1\le
k\le K$, известны из решения на предыдущей $n$-ной итерации.

При переходе от глобальной системы координат $x_l $ к локальной
системе $x_l^{(k)}$ имеют место преобразования тензоров второго
ранга \cite{yank:bib7}
\begin{gather}
\label{yank:eq13} \bar {\sigma }_{ij}^{(k)} =\sum\limits_{l=1}^3
\sum\limits_{m=1}^3 l_{il}^{(k)} l_{jm}^{(k)} \sigma
_{lm}^{(k)},\, \bar {\xi}_{ij}^{(k)} =\sum\limits_{l=1}^3
\sum\limits_{m=1}^3 l_{il}^{(k)} l_{jm}^{(k)} \xi
_{lm}^{(k)},\\ 
i, j=1, 2, 3,\; 1\le k\le K,
\end{gather}
которые с учётом \eqref{yank:eq7} можно записать в матричной
форме:
\begin{equation}
\label{yank:eq14}
\begin{array}{rl}
\displaystyle \bar {\sigma }_k =G_k \sigma _k  & 
\displaystyle  \Bigl(\bar {\sigma }_i^{(k)} =\sum\limits_{j=1}^6 g_{ij}^{(k)} \sigma _j^{(k)}\Bigr),\\
\displaystyle  \bar {\xi}_k = G_k \xi _k, & 
\displaystyle  \Bigl(\bar {\xi}_i^{(k)} =\sum\limits_{j=1}^6 g_{ij}^{(k)} \xi _j^{(k)} \Bigr),\\ 
& i=1, 2, \dots 6.
\end{array}
\end{equation}
Обратные к \eqref{yank:eq14} преобразования имеют вид
\begin{equation}
\label{yank:eq15}
\sigma _k =G_k^{-1} \bar {\sigma }_k,\quad \xi _k =G_k^{-1} \bar {\xi }_k ,\quad 1\le
k\le K,
\end{equation}
где $G_k^{-1} $ "--- матрица, обратная невырожденной ($6{\times}6$)-матрице $G_k =\bigl( g_{ij}^{(k)} \bigr)$ с компонентами
\begin{equation}
\label{yank:eq16}
\begin{array}{l}
g_{11}^{(k)} =l_{11}^{(k)} l_{11}^{(k)},\\
g_{12}^{(k)} =l_{12}^{(k)} l_{12}^{(k)}, \hfill \ldots, \\ 
 g_{66}^{(k)} =l_{11}^{(k)} l_{22}^{(k)} +l_{12}^{(k)} l_{21}^{(k)},\\
\hfill 1\le k\le K;
 \end{array}
\end{equation}
остальные элементы $g_{ij}^{(k)} $ матрицы $G_k$ согласно \eqref{yank:eq13} определяются
табл.~(21.40) в \cite{yank:bib7}, в которой всем величинам нужно приписать верхний индекс
$k$ в скобках, как это сделано в (\ref{yank:eq16}); $l_{ij}^{(k)} $ "--- направляющие косинусы
(см.~(\ref{yank:eq10})); матрица $G_k^{-1}$ совпадает с транспонированной матрицей
(21.44) в \cite{yank:bib7}.

В силу допущений 2, 7 и условий сопряжения полей напряжений, скоростей
установившейся ползучести и температур на границах контакта арматуры со
связующим имеем
\begin{gather}
\label{yank:eq17}
\sum\limits_{j=1}^3 \bar {\sigma }_{ji}^{(k)}
\bar {n}_j^{(k)} =\sum\limits_{j=1}^3 \bar {\sigma
}_{ji}^{(0,k)} \bar {n}_j^{(k)}, \quad i=1,\,2,\,3,\; 1\le k\le K;\\
\label{yank:eq18}
\bar {\xi }_{11}^{(k)} =\bar {\xi }_{11}^{(0,k)} \Rightarrow \bar {\xi }_{11}^{(k)}
=\sum\limits_{l=1}^3 \sum\limits_{m=1}^3 l_{1l}^{(k)}
l_{1m}^{(k)} \xi _{lm}^{(0)},\quad 1\le k\le K;\\
\label{yank:eq19}
\Theta _0 =\Theta _k, \quad 1\le k\le K,
\end{gather}
где $\bar {n}_j^{(k)} $ "--- компоненты единичного вектора ${\rm {\bf n}}_k $,
нормального к поверхности контакта волокна $k$-того семейства со связующим,
заданные в локальной системе координат $x_j^{(k)}$; $\bar {\sigma
}_{ij}^{(0,k)} $, $\bar {\xi }_{11}^{(0,k)} $ "--- компоненты тензоров
напряжений и скоростей деформаций ползучести в связующем, определенные в
локальной системе координат $x_l^{(k)}$.

Так как волокна предполагаются постоянного поперечного сечения (что, как
правило, и имеет место на практике), то в \eqref{yank:eq17} $\bar {n}_1^{(k)} =0$,
поэтому, согласно допущению 3, соотношения \eqref{yank:eq17} будут тождественно
выполняться при произвольных $\bar {n}_2^{(k)}$, $\bar {n}_3^{(k)}$
(таких, что $\bar {n}_2^{(k)2}+\bar {n}_3^{(k)2}=1)$, если выполняются
следующие пять независимых равенств:
\begin{equation}
\label{yank:eq20}
\begin{array}{c}
\displaystyle 
\bar {\sigma }_{ij}^{(k)} =\bar {\sigma }_{ij}^{(0,k)}
\Rightarrow \bar {\sigma }_{ij}^{(k)}=\sum\limits_{l=1}^3 \sum\limits_{m=1}^3 l_{il}^{(k)} l_{jm}^{(k)} \sigma
_{lm}^{(0)},  \\
1\le k\le K;\quad i,\,j=1,\,2,\,3,\quad \text{кроме}\quad i=j=1.
 \end{array}
\end{equation}

Соотношения \eqref{yank:eq17}, \eqref{yank:eq20} определяют статические условия контакта арматуры со
связующим, а равенство \eqref{yank:eq18} "--- кинематическое условие контакта, которое
означает, что на поверхности контакта линейные скорости деформаций
установившейся ползучести в направлении траектории армирования в волокне и в
связующем совпадают, что справедливо в силу допущения 7.

Равенства \eqref{yank:eq18}, \eqref{yank:eq20} перепишем с учётом соотношений \eqref{yank:eq14} и~соответствий~\eqref{yank:eq7}:
\begin{gather}
\label{yank:eq21}
\bar {\xi }_1^{(k)} =\sum\limits_{j=1}^6
g_{1j}^{(k)} \xi _j^{(0)},\quad 1\le k\le K;\\
\label{yank:eq22}
 \bar {\sigma }_i^{(k)} =\sum\limits_{j=1}^6
g_{ij}^{(k)} \sigma _j^{(0)},\quad i=2, 3, \dots, 6,\; 1\le k\le K.
\end{gather}

В силу допущения 3 соотношения \eqref{yank:eq21}, \eqref{yank:eq22} справедливы во всех точках
представительного элемента.

Согласно допущениям 3, 4 усреднённые поля температуры $\Theta $, напряжений
$\sigma $ и скоростей деформаций ползучести $\xi $ в композиции определяются
так:
\begin{gather}
\label{yank:eq23} \Theta =\omega _0 \Theta _0 +\sum\limits_k
\omega _k \Theta _k;\\
\label{yank:eq24} \sigma =\omega _0 \sigma _0 +\sum\limits_k
\omega _k \sigma _k;\\
\label{yank:eq25} \xi =\omega _0 \xi _0 +\sum\limits_k \omega _k
\xi _k.
\end{gather}
Соотношения \eqref{yank:eq24}, \eqref{yank:eq25} записаны в глобальной системе координат $x_l$.

Из равенства \eqref{yank:eq23} с учётом соотношений \eqref{yank:eq19}, \eqref{yank:eq1} и допущения 3 следует
\begin{equation}
\label{yank:eq26}
\Theta =\Theta _k =\Theta _0, \quad 1\le k\le K,
\end{equation}
т.\,е. средняя температура композита в пределах представительного элемента
равна температуре каждого компонента композиции.

%%%

Из равенства \eqref{yank:eq21} с учётом \eqref{yank:eq6}, \eqref{yank:eq26} и допущения 5 вытекает
\begin{equation}
\label{yank:eq27}
\Phi'_k ( S_k ,\Theta ) \sum\limits_{j=1}^6 \bar {a}_{1j}^{(k)} \bar {\sigma }_j^{(k)} = \Phi'_0 (S_0,\,\Theta) \sum\limits_{j=1}^6 g_{1j}^{(k)} \sum\limits_{m=1}^6 a_{jm}^{(0)} \sigma _m^{(0)},\quad 1\le k\le K.
\end{equation}

Систему шести равенств \eqref{yank:eq22}, \eqref{yank:eq27} запишем в матричной форме:
\begin{equation}
\label{yank:eq28}
B_k \bar {\sigma }_k =C_k \sigma _0,\quad 1\le k\le K,
\end{equation}
где ненулевые компоненты ($6{\times}6$)-матриц $B_k =\bigl( b_{ij}^{(k)}\bigr)$, $C_k =\bigl(c_{ij}^{(k)} \bigr)$ согласно \eqref{yank:eq22}, \eqref{yank:eq27} определяются так:
\begin{equation}
\label{yank:eq29}
\begin{array}{ll}
 b_{1j}^{(k)} = \Phi'_k (S_k,\,\Theta)\bar {a}_{1j}^{(k)}, & \displaystyle c_{1j}^{(k)} = \Phi'_0 (S_0,\,\Theta)  \sum\limits_{m=1}^6 g_{1m}^{(k)} a_{mj}^{(0)}, \\
 b_{ii}^{(k)} =1,\quad c_{ij}^{(k)} =g_{ij}^{(k)}, & i=2, 3, \dots, 6, \; j=1, 2, \dots, 6,\; 1\le k\le K.
 \end{array}
\end{equation}
Значения функций $\Phi'_0 (S_0,\,\Theta)$, $\Phi'_k
(S_k,\,\Theta)$ согласно (\ref{yank:eq11}) предполагаются уже
известными из решения на предыдущей $n$-ной итерации.

В силу соотношений (\ref{yank:eq29}) при ненулевом напряжённом состоянии в арматуре $\det B_k \ne 0$, поэтому из \eqref{yank:eq28} получаем
\begin{equation}
\label{yank:eq30}
\bar {\sigma }_k =E_k \sigma _0,\qquad 1\le k\le K,
\end{equation}
где $E_k$ "---  ($6{\times}6$)-матрица:
\begin{equation}
\label{yank:eq31}
 E_k =B_k^{-1} C_k;
\end{equation}
$B_k^{-1}$ "--- ($6{\times}6$)-матрица, обратная $B_k $.

Соотношение \eqref{yank:eq30} определяет на ($n+1)$-й итерации напряжения в $k$-том семействе арматуры $\bar {\sigma }_k $ через напряжения в связующем $\sigma _0 $.

Подставим равенства \eqref{yank:eq15} в \eqref{yank:eq24}, \eqref{yank:eq25} и учтём \eqref{yank:eq11}, \eqref{yank:eq30}, \eqref{yank:eq26}, тогда будем иметь
\begin{equation}
\label{yank:eq32}
\sigma =\Bigl( \omega _0 I +\sum\limits_k \omega _k G_k^{-1} E_k \Bigr)\sigma_0;
\end{equation}
\begin{multline}
\label{yank:eq33}
\xi =2\omega _0 \Phi'_0 (S_0,\,\Theta) A_0 \sigma _0 + 2\sum\limits_k \omega _k \Phi'_k  (S_k,\,\Theta) G_k^{-1} \bar
{A}_k \bar {\sigma }_k = \\
 =2\Bigl[ \omega _0 \Phi'_0 (S_0,\,\Theta)A_0
+\sum\limits_k \omega _k \Phi'_k (S_k ,\,\Theta ) G_k^{-1} \bar {A}_k E_k \Bigr] \sigma _0,
\end{multline}
где $I $ "---  единичная ($6{\times}6$)-матрица.

Из равенства \eqref{yank:eq32} следует
\begin{equation}
\label{yank:eq34} \sigma _0 =H\sigma ,
\end{equation}
где $H$ "--- ($6{\times}6$)-матрица:
\begin{equation}
\label{yank:eq35}
H=\Bigl( \omega _0 I + \sum\limits_k \omega _k G_k^{-1} E_k \Bigr)^{-1}.
\end{equation}

Согласно \eqref{yank:eq29}, \eqref{yank:eq31}, \eqref{yank:eq35} матрица $H$ известна из решения на предыдущей $n$-ной
итерации, поэтому равенство \eqref{yank:eq34} определяет на $(n+1)$-й итерации напряжения в связующем $\sigma _0$ через усреднённые напряжения в композиции~$\sigma$.

Подставим выражение \eqref{yank:eq34} в соотношение \eqref{yank:eq33} и восстановим верхние индексы
$n$ и $n+1$, тогда окончательно получим матричное равенство
\begin{equation}
\label{yank:eq36}
\mathop \xi \limits^{n+1} =2 \Bigl[ \omega_0 \Phi'_0 \bigl(\mathop {S_0 } \limits^n,\,\Theta \bigr)A_0 + 
\sum\limits_k \omega _k \Phi'_k \bigl(\mathop {S_k } \limits^n \Theta  \bigr)G_k^{-1} \bar {A}_k
\mathop {E_k }\limits^n \Bigr] \mathop H\limits^n \mathop
\sigma \limits^{n+1}.
\end{equation}
Из сравнения равенств (\ref{yank:eq12}) и \eqref{yank:eq36} получаем выражение для искомой матрицы:
\begin{equation}
\label{yank:eq37}
\mathop A\limits^n =2 \Bigl[ \omega _0 \Phi'_0 \bigl( \mathop {S_0 }\limits^n,\,\Theta \bigr)A_0
+\sum\limits_k \omega _k \Phi'_k \bigl(\mathop {S_k } \limits^n,\,\Theta \bigr) G_k^{-1} \bar {A}_k \mathop
{E_k }\limits^n \Bigr] \mathop H\limits^n,
\end{equation}
где нужно учесть выражения для ($6{\times}6$)-матриц \eqref{yank:eq35}, \eqref{yank:eq31}, \eqref{yank:eq29}, \eqref{yank:eq16},
\eqref{yank:eq8}, \eqref{yank:eq9} и для направляющих косинусов \eqref{yank:eq10}.

Покажем, как метод последовательных приближений может быть реализован в~рамках предложенной модели механического поведения
пространственно"=армированного металлокомпозита, работающего в условиях установившейся анизотропной ползучести. 
Пусть на некоторой $n$-ной итерации известны приближения напряжений $\mathop {\sigma _0}\limits^n $, $\mathop
{\sigma _k }\limits^n, $ $1\le k\le K$, во всех компонентах композиции, тогда согласно \eqref{yank:eq37}, \eqref{yank:eq35}, \eqref{yank:eq31}, \eqref{yank:eq29}, \eqref{yank:eq8}, \eqref{yank:eq9}, \eqref{yank:eq26} известны компоненты
матрицы $\mathop A\limits^n $ на этой итерации. 
Пусть, кроме того, из решения соответствующей граничной задачи с учётом линейных определяющих
соотношений (\ref{yank:eq12}), где матрица $\mathop A\limits^n $ уже известна, найдены
($n+1)$-е приближения средних напряжений в композиции $\mathop \sigma \limits^{n+1} $. 
Тогда из равенств \eqref{yank:eq34}, \eqref{yank:eq30} с~учётом \eqref{yank:eq35}, \eqref{yank:eq31}
последовательно можем определить следующие приближения напряжений $\mathop
{\sigma _0 }\limits^{n+1} $, $\mathop {\sigma _k }\limits^{n+1} $, $1\le k\le K$, во всех компонентах композиции, после чего на ($n+1)$-й итерации во всех фазовых материалах можем определить значения квадратичных форм $\mathop {S_0 }\limits^{n+1} $, $\mathop {S_k }\limits^{n+1} $ (см.~(\ref{yank:eq2}) с~учётом \eqref{yank:eq7}, \eqref{yank:eq8}). Зная на ($n+1)$-й итерации $\mathop {S_0 }\limits^{n+1} $, $\mathop {S_k }\limits^{n+1} $, $1\le k\le K$, можем определить новые приближения для компонентов матриц $\mathop {B_k }\limits^{n+1} $, $\mathop {C_k }\limits^{n+1} $, $\mathop {E_k }\limits^{n+1} $ (см.~\eqref{yank:eq29}, \eqref{yank:eq31}) и повторить весь алгоритм для получения следующего ($n+2)$-ого приближения
решения и~т.\,д., пока итерационный процесс не сойдётся с требуемой точностью.

Таким образом, соотношения \eqref{yank:eq12} с~учётом равенств \eqref{yank:eq37}, \eqref{yank:eq35}, \eqref{yank:eq31}, \eqref{yank:eq29} характеризуют в матричной форме определяющие уравнения установившейся
анизотропной ползучести пространственно"=армированного металлокомпозита при
реализации метода последовательных приближений.

В общем случае, если деформации ползучести сопровождаются неупругой
дилатацией, т.\,е. $\xi _{11}^{(k)} +\xi _{22}^{(k)} +\xi _{33}^{(k)} =\xi
_1^{(k)} +\xi _2^{(k)} +\xi _3^{(k)} \ne 0$, $0\le k\le K$ (не выполняются
равенства \eqref{yank:eq5}), определяющие соотношения \eqref{yank:eq12} на каждой итерации можно
однозначно разрешить относительно напряжений $\mathop \sigma \limits^{n+1}
=\mathop {A^{-1}}\limits^n \mathop \xi \limits^{n+1} $. Используя при этом
соответствия типа \eqref{yank:eq7} и дифференциальные соотношения Коши \cite{yank:bib3,yank:bib8}
\[
\mathop {\xi _{ij} }\limits^{n+1} =\frac{1}{2}\Bigl( \frac{\partial \mathop
{v_i }\limits^{n+1} }{\partial x_j }+\frac{\partial \mathop {v_j
}\limits^{n+1} }{\partial x_i } \Bigr),\quad i,\,j=1,\,2,\,3,
\]
где $\mathop {v_i }\limits^{n+1} $ "--- $(n+1)$-е приближение скорости установившейся ползучести произвольной точки эквивалентной среды в~направлении $x_i$, можно ставить краевые задачи для металлокомпозитного тела в кинематических переменных "--- скоростях ползучести $v_i$, $i=1,\,2,\,3$.

Однако из экспериментов известно, что дилатация неупругих деформаций, как
правило, пренебрежимо мала \cite{yank:bib3,yank:bib8,yank:bib9}, поэтому при изучении установившейся
ползучести металлов обычно принимают условие несжимаемости \cite{yank:bib3,yank:bib8}
\begin{equation}
\label{yank:eq38}
\xi _{11}^{(k)} +\xi _{22}^{(k)} +\xi
_{33}^{(k)} =\xi _1^{(k)} +\xi _2^{(k)} +\xi _3^{(k)}
=0, \quad 0\le k\le K.
\end{equation}

Согласно \eqref{yank:eq1}, \eqref{yank:eq25}, \eqref{yank:eq38} и допущениям 3, 4, несжимаемостью обладает и~материал композиции:
\[
\mathop {\xi _{11} }\limits^{n+1} +\mathop {\xi _{22}}\limits^{n+1}
+\mathop {\xi _{33} }\limits^{n+1} =\mathop {\xi _1 }\limits^{n+1} +\mathop
{\xi _2 }\limits^{n+1} +\mathop {\xi _3 }\limits^{n+1} =0,
\]
откуда с учётом соответствий типа \eqref{yank:eq7} вытекает линейная зависимость первых трёх уравнений в~системе \eqref{yank:eq12}. 
Следовательно, в случае \eqref{yank:eq38} (при выполнении равенств \eqref{yank:eq5}, \eqref{yank:eq9}) систему \eqref{yank:eq12} нельзя на каждой итерации однозначно разрешить относительно $(n+1)$-х приближений напряжений в~композиции 
$\mathop \sigma \limits^{n+1} $, а значит, при выполнении равенств \eqref{yank:eq38},
\eqref{yank:eq5}, \eqref{yank:eq9} предложенная модель механического поведения пространственно"=армированного металлокомпозита, работающего в условиях установившейся ползучести, при выполнении допущения~7 может быть использована на практике лишь в тех задачах, которые допускают решение в напряжениях $\sigma _{ij} $, а не в кинематических переменных "--- скоростях
ползучести $v_i $. Однако в реальности большинство конструкций имеют такие типы закрепления и нагружения, что соответствующие граничные задачи для них нужно решать в кинематических переменных, и предложенная модель в этих случаях малопригодна, когда имеют место равенства \eqref{yank:eq38}, \eqref{yank:eq5}, \eqref{yank:eq9}.

Если отказаться от выполнения допущения 7 и учитывать упругие (или упругопластические) деформации в компонентах композиции, то придётся  рассматривать, по сути, неустановившуюся ползучесть этих материалов \cite{yank:bib3} и композиции в целом. При этом вместо равенства \eqref{yank:eq18} нужно использовать
соотношение
\begin{multline}
\label{yank:eq39}
\bar {\xi }_{11}^{(k)} +\dot {\bar {e}}_{11}^{(k)} =\bar {\xi
}_{11}^{(0,k)} +\dot {\bar {e}}_{11}^{(0,k)} \Rightarrow \\
 \Rightarrow  \bar {\xi }_{11}^{(k)} +\dot {\bar {e}}_{11}^{(k)}
=\sum\limits_{l=1}^3 \sum\limits_{m=1}^3 l_{1l}^{(k)} l_{1m}^{(k)} \left(\xi _{lm}^{(0)} +\dot {e}_{lm}^{(0)}\right),\quad 1\le k\le K,
\end{multline}
где $\dot {e}_{lm}^{(0)}$, $\dot {\bar {e}}_{11}^{(k)} $ "--- скорости
упругих деформаций в компонентах композиции; \linebreak $\dot {\bar {e}}_{11}^{(0,k)}$
"--- скорость упругих деформаций в связующем в направлении $x_1^{(k)} $ "--- в~направлении траектории армирования волокнами $k$-того семейства.

Используя закон Гука, в \eqref{yank:eq39} следует выразить $\dot {\bar {e}}_{11}^{(k)} $,
$\dot {e}_{lm}^{(0)} $ через скорости напряжений $\dot {\bar {\sigma
}}_{ij}^{(k)}$, $\dot {\sigma }_{ij}^{(0)} $ соответственно, после чего
получим зависимость, содержащую в явном виде производные по времени от $\bar
{\sigma }_{ij}^{(k)} $, $\sigma _{ij}^{(0)} $, поэтому предложенный выше
подход к моделированию механического поведения пространственно"=армированного
металлокомпозита в условиях неустановившейся анизотропной ползучести
становится не~приемлемым. Моделирование же неустановившейся ползучести
металлокомпозитов выходит за рамки настоящего исследования.

\smallskip

\Section{Сравнительный анализ моделей установившейся ползучести
перекрестно армированных металлокомпозитов} Автору неизвестны
экспериментальные данные по установившейся ползучести металлокомпозитов
перекрестной структуры, поэтому верификацию предложенной модели можно
провести лишь косвенно. Известно, что в изотропном случае определяющие
квазилинейные уравнения теории установившейся ползучести, использованные в
п.~1, и деформационной теории пластичности формально полностью
совпадают \cite{yank:bib3,yank:bib8,yank:bib9} (если вместо скоростей деформаций ползучести понимать
просто деформации). Частным же случаем определяющих уравнений деформационной
теории является закон Гука. Поэтому если при моделировании упругого
поведения перекрестно армированного композита на основе допущений,
аналогичных использованным в настоящем исследовании, удастся получить
расчётные значения эффективных упругих характеристик композиции,
удовлетворительно согласующиеся с экспериментом, то тем самым получим
косвенное подтверждение приемлемости исходных гипотез данного исследования и
модели, построенной на их основе. В работах [10--12] автором были предложены
модели упругого поведения однонаправленно и перекрестно армированных
гибридных 2D- и 3D-композитов, основанные на исходных предпосылках,
аналогичных использованным в настоящем исследовании. В работах же \cite{yank:bib11,yank:bib13,yank:bib14}
продемонстрировано, что расчётные значения эффективных упругих
характеристик, определенные по этим моделям, удовлетворительно согласуются с
экспериментальными данными. Это обстоятельство позволяет доверительно
относиться к исходным допущениям, использованным в данном исследовании, и к
модели, построенной на их основе. Окончательный же вывод о приемлемости этой
модели может быть сделан лишь после проведения серии соответствующих
экспериментальных испытаний об установившейся ползучести перекрестно
армированных металлокомпозитов.

В настоящем же пункте проведём сравнение механического поведения
перекрестно армированных в плоскости (2D) и в пространстве (3D)
металлокомпозитов в условиях установившейся ползучести, рассчитанного
(поведения) на основе разных структурных моделей: 1)~модели, предложенной в
п.~1; 2)~модели Ю.\,В.~Немировского с одномерным напряжённым состоянием в~волокнах \cite{yank:bib6}.

С этой целью рассмотрим два образца. Первый образец"=полоска выполнен из меди
(Cu) и ортогонально армирован в плоскости ($x_1$, $x_2$) двумя
семействами стальной проволоки У8А ($K=2$) под углами $\varphi _1
=\rm{const}$, $\varphi _2 =\varphi _1 -\pi /2=\rm{const}$ ($\theta _1
=\theta _2 =\pi /2$, см. (\ref{yank:eq10})) с плотностями $\omega _1 =\omega _2 =0,128$.
Образец нагружен растягивающим напряжением $\sigma _{11} =1$~ГПа,
остальные напряжения равны нулю (рис.~2). Образцы"=полоски такого типа, как
правило, используются при испытаниях армированных композитов на растяжение
\cite{yank:bib15}.

\begin{figure}[b!]\centering
\includegraphics[scale=.9]{./00PIC/yank2.eps}
\caption{Расчётная схема армированного образца}
\end{figure}

Второй 3D-образец также выполнен из меди и пространственно армирован в трёх
($K=3$) ортогональных направлениях проволокой У8А с плотностями $\omega _1
=\omega _2 =0,128$, $\omega _3 =0,333$ (при этом суммарная плотность
армирования \mbox{$\omega _1 +\omega _2 +\omega _3 =0,589$} соответствует структуре
с плотной упаковкой, см.~табл.~1.2 в~\cite{yank:bib1}). Направления армирования
определяются углами сферической системы координат (см.~(\ref{yank:eq10})): $\theta _1
=\pi /2$, $\varphi _1 =\rm{const}$, $\theta _2 =\pi /2$, $\varphi _2
=\varphi _1 -\pi /2=\rm{const}$, $\theta _3 =\varphi _3 =0$, т.\,е. два
первых семейства проволок У8А уложены так, как изображено на рис.~2, а
третье семейство армирования перпендикулярно плоскости этого рисунка. Второй
образец нагружен так же, как и первый.

Образцы рассчитываются при температуре $\Theta \approx 200~^\circ$С, при которой характеристики степенного закона установившейся
ползучести $\xi _k =b_k \sigma ^{m_k }$ ($0\le k\le K$) для материалов
компонентов композиции, которые предполагаются изотропными, имеют значения:
\begin{equation}
\label{yank:eq40}
\begin{array}{rll}
 \rm Cu \quad [16]:    & m_0 =2,16,  & b_0 =5,63 \cdot 10^{-9}~(\text{МПа})^{-m_0}\cdot \text{ч}^{-1}, \\
 \text{У8А}\quad [15]: & m_1 =24,98, & b_1 =1,054 \cdot 10^{-84}~(\text{МПа})^{-m_1 }\cdot \text{ч}^{-1}.
 \end{array}
\end{equation}
(В изотропном случае квадратичные формы (\ref{yank:eq2}) с учётом (\ref{yank:eq3}), (\ref{yank:eq8}), (\ref{yank:eq9}) при
$\alpha _k =\beta _k =0$ пропорциональны второму инварианту тензора
напряжений в $k$-том компоненте композиции, а потенциалы ползучести $\Phi_k
(S_k,\,\Theta _k)$, $0\le k\le K$, определяются
соотношениями (68.1) в \cite{yank:bib3} с учётом (\ref{yank:eq26}), \eqref{yank:eq40}.)

На рис.~3,~4 изображены зависимости усредненных скоростей деформаций
установившейся ползучести $\xi _{ij}$ рассматриваемых металлокомпозитных
2D- (рис.~3) и 3D-образцов (рис.~4) от угла армирования первым семейством волокон $\varphi _1 $ (см.~рис.~2). 
На этих рисунках кривые~1 характеризуют зависимости $\xi _{11} ( \varphi _1 )$, линии 2 "--- $\xi _{22}(\varphi _1)$,
кривые 3 "--- $\xi _{33}(\varphi _1)$, линии 4 "--- $\xi _{12} (\varphi _1)$ ($\xi _{13}
\equiv \xi _{23} \equiv 0$). 
Кривые 5 определяют зависимость интенсивности скоростей деформаций сдвига $H$ (см.~(12.11) в \cite{yank:bib8}) в композиции от угла армирования $\varphi _1$. На рис.~3, $a$ и 4, $a$ представлены результаты
расчётов по структурной модели, предложенной в п.~1, а на~рис.~3,~{\sl б} и 4,~{\sl б} "--- по модели 
Ю.\,В.~Немировского~\cite{yank:bib6}.

\begin{figure}[b!]\centering \small
\includegraphics[width=\textwidth]{./00PIC/yank3.eps}\\
{\sl а \hspace{60mm} б}
\caption{Зависимости скоростей деформаций установившейся ползучести 2D"=металлокомпозитного образца от угла армирования
первым семейством: {\sl а} "--- расчёт по структурной модели из п.~1; {\sl б} "--- расчёт по структурной модели 
Ю.\,В.~Немировского
\cite{yank:bib6,yank:bib17}}
%\end{figure}

\bigskip

%\begin{figure}[h]\centering
\includegraphics[width=\textwidth]{./00PIC/yank4.eps}\\
{\sl а \hspace{60mm} б}
\caption{Зависимости скоростей деформаций установившейся ползучести 3D"=металлокомпозитного образца от угла армирования
первым семейством: {\sl а} "--- расчёт по структурной модели из п.~1; {\sl б} "--- расчёт по структурной модели 
Ю.\,В.~Немировского
\cite{yank:bib6,yank:bib18}}
\end{figure}

Структурные соотношения установившейся ползучести перекрестно армированного
слоя 2D-металлокомпозита, соответствующие модели Ю.\,В.~Не\-мировского при
плоском напряжённом состоянии, приведены в \cite{yank:bib17} (см.~там формулы (\ref{yank:eq3}), (\ref{yank:eq4}),
(\ref{yank:eq7})). Эти структурные соотношения элементарно обобщаются на случай
пространственного армирования 3D-композита \cite{yank:bib18}. Следует лишь использовать
определяющие соотношения для трехмерного напряжённого состояния в связующем,
а под направляющими косинусами $l_{ik}$ в формуле (\ref{yank:eq3}) из \cite{yank:bib17} понимать
$l_{1i}^{(k)}$, $i=1,\,2,\,3$, $1\le k\le K$, принятые в настоящем
исследовании (см.~(\ref{yank:eq10})).

Из шести условий сопряжения полей перемещений и напряжений на границах
контакта волокон со связующим (\ref{yank:eq21}), (\ref{yank:eq22}) (тождественных (\ref{yank:eq18}), (\ref{yank:eq20})) в модели
из \cite{yank:bib6,yank:bib17,yank:bib18} используется только кинематическое соотношение (\ref{yank:eq21}) (или, что
то же самое, (\ref{yank:eq18})). Статические же условия контакта (\ref{yank:eq22}) вообще не
используются, поэтому в рамках этой модели получаются достаточно простые
структурные соотношения, но, как уже отмечалось во введении, напряжённое
состояние в волокнах в поперечном направлении при этом не учитывается, а
значит невозможно учесть и анизотропию ползучести в арматуре, так как при
этом в (\ref{yank:eq2}) все напряжения кроме $\sigma _{11}^{(k)}$, $1\le k\le K$,
принимаются тождественно равными нулю. В силу указанных обстоятельств при
некоторых углах рассогласования направлений армирования и нагружения
усиливающие волокна как бы не участвуют в работе композита при его
деформировании, и в этих случаях большую часть нагрузки принимает на себя
связующее, которое при этом как бы ослаблено наличием неработающей (или
частично работающей) арматуры. Результатом этого явился тот факт, что
максимальные по модулю значения ординат точек кривых на рис.~3,~{\sl б} почти вдвое
больше, чем на рис.~3,~{\sl а}; на рис.~4, {\sl б} "--- почти в шесть раз больше, чем на
рис.~4, {\sl а}. (Качественное поведение кривых с одинаковыми номерами на рис.~3,~4
подобно, различие заключается лишь в значениях ординат точек на этих
кривых.)

В рамках же модели, предложенной в п.~1 настоящего исследования,
учитываются все условия сопряжения (\ref{yank:eq21}), (\ref{yank:eq22}) на поверхностях контакта
арматуры со связующим, что соответствует идеальной адгезии и приводит к
меньшей расчётной податливости металлокомпозита в условиях установившейся
ползучести по сравнению с расчётом по модели Ю.\,В.~Немировского \cite{yank:bib6,yank:bib17,yank:bib18}.

Расчёт изотропного образца из меди согласно (\ref{yank:eq40}) приводит к значению $\xi
_{11} =b_0 \sigma _{11}^{m_0 } =17\cdot 10^{-3}$~ч$^{-1}$, что на
несколько порядков больше значений ординат точек на кривых $1$ рис.~3,~4.
Следовательно, в рассматриваемом случае обе структурные модели
предсказывают, что армирование образца проволокой У8А приводит к
существенному снижению его податливости в условиях установившейся
ползучести, что оправдывает целесообразность армирования.

Поведение всех кривых на рис.~3,~4 показывает, что изменение структуры
армирования существенно влияет на механическое поведение металлокомпозита в
условиях установившейся ползучести. В частности, скорости деформаций
ползучести $\xi _{ij} $ металлокомпозитного образца при разных углах
армирования $\varphi _1 $ могут различаться на несколько порядков. Поэтому
имеет смысл осуществлять целевую оптимизацию армированных конструкций,
работающих в условиях установившейся ползучести.

В силу указанных выше особенностей сравниваемых структурных моделей
оптимальные проекты армирования, полученные на их основе, могут и не
совпадать. Так, сопоставление ординат точек на кривых рис.~3,~{\sl а} и 4,~{\sl а}
показывает, что дополнительное армирование третьим семейством волокон в
направлении $x_3 $ приводит к уменьшению податливости металлокомпозита почти
вдвое в рамках расчёта по модели, предложенной в п.~1; аналогичное же
сравнение кривых на рис.~3, {\sl б} и 4,~{\sl б} демонстрирует, что в рамках модели Ю.\,В.~Немировского
\cite{yank:bib6,yank:bib18} дополнительное армирование в направлении $x_3$,
наоборот, вызывает увеличение податливости металлокомпозитного образца тоже
почти вдвое. Последнее обстоятельство объясняется тем, что, согласно модели
из \cite{yank:bib6,yank:bib18}, третье семейство арматуры не воспринимает нагружение в плоскости
($x_1 ,\,x_2 $), а лишь ослабляет связующее.

Следовательно, если в качестве критерия эффективного армирования выступает
требование уменьшения скоростей деформаций ползучести, то, согласно модели
из п.~1, использование третьего семейства арматуры повышает
эффективность работы образца, а согласно модели Ю.\,В.~Немировского, наоборот,
"--- уменьшает.

\smallskip

\Section[a]{Заключение} Предложенная структурная модель позволяет строить в
итерационной форме определяющие уравнения для пространственно"=армированных в
произвольных направлениях металлокомпозитов при наличии идеальной адгезии
между связующим и арматурой, материалы которых работают в условиях
установившейся анизотропной ползучести. Зная же начальное механическое
состояние металлокомпозита \cite{yank:bib12} и его состояние в условиях установившейся
ползучести и используя экстремальные принципы теории ползучести \cite{yank:bib8}, можно
получить приближенное решение о неустановившейся ползучести такого композита
на первой её стадии.

\thanks{Работа выполнена при поддержке РФФИ (код проекта 12--01--90405--Укр\_а).}



\begin{thebibliography}{99}

\RBibitem{yank:bib1}
\book Пространственно"=армированные композиционные материалы
\bookinfo Справочник
\by Тарнопольский~Ю.\,М., Жигун~И.\,Г., Поляков~В.\,А.
\publaddr М.
\publ Машиностроение
\yr 1987
\totalpages 224
\misctransl
\by  Tarnopolsky~Yu.\,M.,  Zhigun~I.\,G.,  Polyakov~V.\,A.
\book Spatial-Reinforced Composite Materials 
\bookinfo Handbook
\publ Mashinostrocnie
\publaddr Moscow 
\yr 1987
\totalpages 224

\Bibitem{yank:bib2}
\by Mohamed~M.\,H., Bogdanovich~A.\,E., Dickinson~L.\,C., Singletary~J.\,N., Lienhart~R.\,R.
\paper A New Generation of 3D Woven Fabric Preforms an Composites
\jour SAMPE J.
\yr 2001
\vol 37
\issue 3
\pages 8--17

\RBibitem{yank:bib3}
\by Работнов~Ю.\,Н.
\book Ползучесть элементов конструкций
\publaddr М.
\publ Наука
\yr 1966
\totalpages 752
\misctransl
\by  Rabotnov~Yu.\,N.
\book Creep of Structural Elements
\publ Nauka
\publaddr Moscow
\yr 1966
\totalpages 752

\RBibitem{yank:bib4}
\by Резников~Б.\,С., Никитенко~А.\,Ф., Кучеренко~И.\,В.
\paper Установившаяся ползучесть микронеоднородных сред
\inbook Труды четвёртой Всероссийской научной конференции с~международным участием \rm (29--31 мая 2007~г.). Часть~1
\bookinfo Математические модели механики, прочности и надёжности элементов конструкций
\serial Матем. моделирование и краев. задачи
\yr 2007
\pages 219--223
\publ СамГТУ
\publaddr Самара
\misctransl
\by Reznikov~B.\,S., Nikitenko~A.\,F., Kucherenko~I.\,V.
\paper Steady creep of micrononuniform media
\inbook Proceedings of the Fourth All-Russian Scientific Conference with international participation \rm (29--31 May 2007). Part~1
\serial Matem. Mod. Kraev. Zadachi
\yr 2007
\pages 219--223
\publ SamGTU
\publaddr Samara


\RBibitem{yank:bib5}
\by Каримбаев~Т.\,Д., Мыктыбеков~Б.\,М., Панова~И.\,М.
\book Математические модели нелинейного деформирования однонаправленно"=армированных композиционных материалов
\serial Труды ЦИАМ
\publaddr М.
\publ ЦИАМ
\yr 2005
\vol 1334
\totalpages 160
\misctransl 
\by Karimbaev~T.\,D., Myktybekov~B.\,M.,  Panova~I.\,M.
\book Mathematical models of nonlinear deformation of unidirectionally reinforced composite materials
\serial Trudy CIAM 
\publaddr Moscow
\publ CIAM
\yr 2005
\vol 1334
\totalpages 160


\RBibitem{yank:bib6}
\by Немировский~Ю.\,В.
\paper Ползучесть однородных и композитных оболочек
\inbook Актуальные проблемы механики оболочек
\bookinfo Тр. междунар. конф., посвящённой 100-летию проф.~Х.~М.~Муштари, 90-летию проф. К.~З.~Галимова
и 80-летию проф. М.~С.~Корнишина 
\procinfo Казань, 26--30 июня 2000~г.
\publaddr Казань
\publ Новое знание
\yr 2000
\pages 42--49
\misctransl
\by Nemirovskii~Yu.\,V.
\paper Creep of homogeneous and composite shells
\inbook Topical Problems of the Mechanics of Shells
\bookinfo Trans. Int. Conf. Dedicated to the 100th Birthday of Prof. Kh.~M.~Mushtari, 90th Birthday of Prof. K.~Z.~Galimov, and 80th Birthday of Prof. M.~S.~Kornishin
\procinfo Kazan', 26--30 June, 2000
\publ Novoe Znanie
\publaddr Kazan'
\yr 2000
\pages 42--49



\RBibitem{yank:bib7}
\by Малмейстер~А.\,К., Тамуж~В.\,П., Тетерс~Г.\,А.
\book Сопротивление жестких полимерных материалов
\publaddr Рига
\publ Зинатне
\yr 1972
\totalpages 500
\misctransl
\by Malmeister~A.\,K., Tamuzh~V.\,P.,  Teters~G.\,A.
\book Strength of Polymer and Composite Materials 
\publ Zinatne
\publaddr Riga
\yr 1972
\totalpages 500

\RBibitem{yank:bib8}
\by Качанов~Л.\,М.
\book Теория ползучести
\publaddr М.
\publ Физматгиз
\yr 1960
\totalpages 456
\misctransl
\by Kachanov~L.\,M.
\book Creep Theory 
\publ Fizmatgiz
\publaddr Moscow 
\yr 1960
\totalpages 456

\RBibitem{yank:bib9}
\by Ильюшин~А.\,А.
\book Труды
\bookvol 3
\voltitle Теория термовязкоупругости
\eds Е.~А.~Ильюшина, В.~Г.~Тунгускова
\publaddr М.
\publ  Физматлит
\yr 2007
\totalpages 288
\misctransl
\by Il'yushin~A.\,A.
\book Works
\eds E.~A.~Il'yushina, V.~G.~Tunguskova
\bookvol 3
\voltitle Theory of Thermoviscoelasticity 
\publ Fizmatlit
\publaddr Moscow 
\yr 2007
\totalpages 288


\RBibitem{yank:bib10}
\by Немировский~Ю.\,В., Янковский~А.\,П.
\paper Эффективные физико"=механические характеристики композитов, однонаправленно армированных монотропными волокнами.
 Сообщение 1. Модель армированной среды
\jour Изв. вузов. Строительство
\yr 2006
\issue 5
\pages 16--24
\misctransl
\by Nemirovsky~Yu.\,N., Yankovsky~A.\,P.
\paper Effective physicomechanical characteristics of composites unidirectionally reinforced with transversely isotropic fibers. Report~1: Model of a reinforced medium
\jour Izv. vuzov. Stroitel'stvo 
\yr 2006
\issue 5
\pages 16--24

\RBibitem{yank:bib11}
\by Немировский~Ю.\,В., Янковский~А.\,П.
\paper Определение эффективных физико"=механических характеристик гибридных композитов, перекрестно армированных трансверсально"=изотропными волокнами, и сопоставление расчётных характеристик с экспериментальными данными
\jour Мех. композ. матер. и констр.
\yr 2007
\vol 43
\issue 1
\pages 3--32
\misctransl
\by Nemirovsky~Yu.\,N., Yankovsky~A.\,P.
\paper Determination of the effective physicomechanical characteristics of hybrid composites cross-reinforced with transversely isotropic fibers and a comparison of calculated characteristics with experimental data
\jour Mekh. Kompozits. Mater. Konstr.
\yr 2007
\vol 43
\issue 1
\pages 3--32


\RBibitem{yank:bib12}
\by Янковский~А.\,П.
\paper Определение термоупругих характеристик пространственно армированных волокнистых сред при
общей анизотропии материалов компонент композиции. 1. Структурная модель
\jour Мех. композ. матер.
\yr 2010
\vol 46
\issue 5
\pages 663--678
\transl англ. пер.:~
\paper Determination of the thermoelastic characteristics of spatially reinforced fibrous media in the case of general anisotropy of their components. 1.~Structural model
\by Yankovskii~A.\,P.
\jour Mech. Compos. Mater.
\vol 46
\issue 5
\yr 2010
\pages 451--460



\RBibitem{yank:bib13}
\by Немировский~Ю.\,В., Янковский~А.\,П.
\paper Эффективные физико"=механические характеристики композитов, однонаправлено армированных монотропными волокнами. Сообщение~2. Сопоставление расчётных характеристик с экспериментальными данными
\jour Изв. вузов. Строительство
\yr 2006
\issue 6
\pages 10--19
\misctransl
\by Nemirovsky~Yu.\,N., Yankovsky~A.\,P.
\paper Effective physicomechanical characteristics of composites unidirectionally reinforced with transversely isotropic fibers. Report~2: Comparison between calculated characteristics and experimental data
\jour Izv. vuzov. Stroitel'stvo 
\yr 2006
\issue 6
\pages 10--19


\RBibitem{yank:bib14}
\by Янковский~А.\,П.
\paper Определение термоупругих характеристик пространственно армированных волокнистых сред при общей анизотропии
 материалов компонент композиции. 2.~Сравнение с~экспериментом
 \jour Мех. композ. матер.
 \yr 2010
 \vol 46
 \issue 6
 \pages 955--964
 \transl англ. пер.:~
\jour Mech. Compos. Mater.
\vol 46
\issue 6
\pages 659--666
%, DOI: 10.1007/s11029-011-9179-9
\paper Determination of the thermoelastic characteristics of spatially reinforced fibrous media in the case of general anisotropy of their components. 2.~Comparison with experiment
\by Yankovskii~A.\,P.


\RBibitem{yank:bib15}
\book Композиционные материалы
\bookinfo Справочник
\ed Д.\,М.~Карпинос
\publaddr Киев
\publ Наукова думка
\yr 1985
\totalpages 592
\misctransl
\ed D.~M.~Karpinos 
\book Composite Materials
\bookinfo Handbook 
\publ Naukova Dumka
\publaddr Kiev 
\yr 1985
\totalpages 592

\RBibitem{yank:bib16}
\by Писаренко~Г.\,С., Можаровский~Н.\,С.
\book Уравнения и краевые задачи теории пластичности и ползучести
\bookinfo Справочное пособие
\publaddr Киев
\publ Наукова думка
\yr 1981
\totalpages 496
%http://www.ams.org/mathscinet-getitem?mr=645541
%http://www.zentralblatt-math.org/zmath/search/?an=Zbl%200511.73029
\misctransl
\by Pisarenko~G.\,S.,  Mozharovskii~N.\,S.
\book Equations and Boundary-Value Problems of the Theory of Plasticity and Creep
\bookinfo Handbook 
\publ Naukova Dumka
\publaddr Kiev 
\yr 1981
\totalpages 496

\RBibitem{yank:bib17}
\by Янковский~А.\,П.
\paper Установившаяся ползучесть сложно армированных металлокомпозитных пластин, нагруженных в своей плоскости
\jour Мат. модел.
\yr 2010
\vol 22
\issue 8
\pages 55--66
%Zbl pre05863682 
\misctransl
\by Yankovskii~A.\,P.
\paper The steady creeping difficulty reinforced the metal-composite plates loaded in the plane 
\jour Mat. Model. 
\vol 22
\issue 8
\pages 55--66 
\yr 2010

\RBibitem{yank:bib18}
\by Немировский~Ю.\,В., Янковский~А.\,П.
\paper Численное моделирование нелинейно"=наследственного поведения пространственно"=армированных композитных сред
\inbook Численные методы решения задач теории упругости и пластичности
\bookinfo Тез. докл. XXII Всерос. конф.
\procinfo Барнаул, 4--7 июля 2011~г.
\publaddr Новосибирск
\publ Параллель
\yr 2011
\pages 53
\misctransl
\by Nemirovsky~Yu.\,N., Yankovsky~A.\,P.
\paper Numerical simulation of behaviour of the three-dimensional reinforcement composit materials with nonlinear memory
\inbook Numerical methods for solving problems of elasticity and plasticity theory 
\publaddr Novosibirsk
\publ Parallel'
\yr 2011
\pages 53

\end{thebibliography}


\noindent
\hskip6.5cm \thedate \par\noindent
\hskip6.5cm\therevdate


\vspace{-6mm}
%
%\chead[\fancyplain{}{\scriptsize \sl Y\,a\,n\,k\,o\,v\,s\,k\,i\,i~A.\,P.}]{\fancyplain{}{\scriptsize \sl  Modeling of the Steady-State Creep of Three-Dimensional
%Reinforced \dots}}

\makeengtitle

\newpage

%
% \tableofengcontents
%%
%\end{document}
